\section{Experiments}
\label{app-experiments}

This section aims at providing the reader with details of the experiments reported in Section~\ref{sect-experiments}.

Table~\ref{tab:fed-partitions} shows the different distributions of positive and negative examples across clients for each dataset. We use this partitioning across all federated experiments. 

\begin{table}[ht]
\centering

\setlength{\tabcolsep}{6pt}
\begin{tabular}{l|ccc|ccc|ccc}
\toprule
\textbf{Partition}
& \multicolumn{3}{c|}{\textbf{Trains}}
& \multicolumn{3}{c|}{\textbf{Zendo1}}
& \multicolumn{3}{c}{\textbf{IGGP-RPS}} \\
& POS & NEG & Total
& POS & NEG & Total
& POS & NEG & Total \\
\midrule

Part 1
& 2 & 2 & 4
& 7 & 7 & 14
& 36 & 119 & 155 \\

Part 2
& 2 & 2 & 4
& 7 & 7 & 14
& 36 & 119 & 155 \\

Part 3
& 1 & 1 & 2
& 6 & 6 & 12
& 36 & 118 & 154 \\

\bottomrule \\
\end{tabular}
\caption{Positive and negative examples per dataset partition}
\label{tab:fed-partitions}
\end{table}

Algorithm~\ref{alg:gsrvpopper} details the algorithm used to conduct the correctness study.

\begin{algorithm}[!t]
  \caption{Srvpopper client for correctness tests}
  \label{alg:gsrvpopper}
  \small
  \begin{lstlisting}[breaklines=true, breakatwhitespace=true]
1  def gsrvpopper(N, E+, E-, B, D, C, 
                       max_vars, max_literals, max_clauses):
2  
3      CC = C
4      best_score = 0
5      best_hyp = { }
6      end_popper = false
7      num_literals = 0
8      r = 1
9  
10     while not end_popper and num_literals <= max_literals:
11        c_solvable = true
12        num_literals += 1
13        while not end_popper and c_solvable:
14           hyp = generate(D, CC, max_vars, num_literals, 
                                                 max_clauses)
15           if not hyp:
16              c_solvable = false
17           else:
18              // Global computation
19              conf_matrix = test(E+, E-, B, hyp)
20              outcome_g = compute_outcome(conf_matrix)
21              score_g = compute_score(conf_matrix)
22             // Federated computation
23              tell(hypsi(r,hyp))
24              for i in range(1,N+1):
25                 cm[i] = ask(cmtuple(r,i))
26              outcome_fed = aggregate_outcome(cm)
27              score_fed = aggregate_score(cm)
28              // Results comparison
29              if (outcome_g != outcome_fed) ||
                     (score_g != score_fed):
30                  print("Global calculus non consistent with 
                             federated calculus!")
31              if outcome == ('all', 'none'):
32                 best_hyp = hyp
33                 end_popper = true
34                 tell(hypsi(r,nil_hyp))                
35              else:
36                 if best_score < score:
37                    best_hyp = hyp
38                    cc += learn_constraints(hyp, outcome)
39            r = r+1
40                    
41     return best_hyp
  \end{lstlisting}
\end{algorithm}




Algorithm~\ref{alg:tgsrvpopper} details the algorithm used to conduct the performance study.

\begin{algorithm}[!t]
  \caption{Srvpopper client for performance tests}
  \label{alg:tgsrvpopper}
  \small
  \begin{lstlisting}[breaklines=true, breakatwhitespace=true]
1  def tgsrvpopper(N, E+, E-, B, D, C, 
                      max_vars, max_literals, max_clauses):
  
2      start_time = time.process_time()     
3      CC = C
4      best_score = 0
5      best_hyp = { }
6      end_popper = false
7      num_literals = 0
8      r = 1
9      tCentralPopper = 0
10     tFedPopper = 0
11  
12     while not end_popper and num_literals <= max_literals:
13        c_solvable = true
14        num_literals += 1
15        while not end_popper and c_solvable:
16           hyp = generate(D, CC, max_vars, 
                                   num_literals, max_clauses)
17           if not hyp:
18              c_solvable = false
19           else:
20              // Global computation
21              start_CentralTime = time.process_time()
22              conf_matrix = test(E+, E-, B, hyp)
23              outcome_g = compute_outcome(conf_matrix)
24              score_g = compute_score(conf_matrix)
25              end_CentralTime = time.process_time()
26              tCentralPopper = tCentralPopper + 
                        (end_CentralTime - start_CentralTime)
27             // Federated computation
28                start_FedTime = time.process_time()
29              tell(hypsi(r,hyp))
30              for i in range(1,N+1):
31                 cm[i] = ask(cmtuple(r,i))
32              outcome_fed = aggregate_outcome(cm)
33              score_fed = aggregate_score(cm)
34              end_FedTime = time.process_time()
35              tFedPopper = tFedPopper + 
                        (end_FedTime - start_FedTime)
36              if outcome == ('all', 'none'):
37                 best_hyp = hyp
38                 end_popper = true
39                 tell(hypsi(r,nil_hyp))                
40              else:
41                 if best_score < score:
42                    best_hyp = hyp
43                    cc += learn_constraints(hyp, outcome)
44            r = r+1
             
45      end_time = time.process_time()
46      global_time = end_time - start_time
47      print("Time for Central Popper : %10d", 
                            %(global_time - tFedPopper))
48      print("Time for Fed Popper : %10d", 
                            %(global_time - tCentralPopper))
49      return best_hyp
  \end{lstlisting}
\end{algorithm}
